\section{Learning Objectives}

After completing this chapter, readers will be able to:
\begin{itemize}
    \item Implement comprehensive data validation strategies
    \item Use Great Expectations for data quality testing
    \item Design data quality metrics and monitoring
    \item Handle data anomalies and outliers
    \item Build automated data quality pipelines
\end{itemize}

\section{Data Quality Dimensions}

\subsection{The Six Dimensions of Data Quality}

\begin{itemize}
    \item \textbf{Accuracy:} Data correctly represents reality
    \item \textbf{Completeness:} All required data is present
    \item \textbf{Consistency:} Data is consistent across systems
    \item \textbf{Timeliness:} Data is up-to-date
    \item \textbf{Validity:} Data conforms to defined formats
    \item \textbf{Uniqueness:} No duplicate records
\end{itemize}

\section{Automated Data Validation}

% Placeholder content
\section{Data Quality Metrics}

% Placeholder content
\section{Handling Data Quality Issues}

% Placeholder content
\section{Chapter Summary}

% Placeholder summary
\section{Exercises}

\begin{exercise}
Implement a comprehensive data validation suite for a customer dataset using Great Expectations.
\end{exercise}
